% ============================================================
% CAPÍTULO 1 — INTRODUCCIÓN (IEEE 830 §1)
% ============================================================
\chapter{Introducción}

Este documento constituye la Especificación de Requerimientos de Software (ERS)
del \textbf{\nombreSistema}, elaborada conforme al estándar \textbf{IEEE 830}.
Todos los requerimientos, reglas de negocio y casos de uso aquí descritos
fueron extraídos directamente del código fuente y de la base de datos del
proyecto, por lo que reflejan con exactitud el comportamiento real del sistema
implementado.

\section{Propósito}

El propósito de este documento es especificar de forma completa, precisa y
verificable los requerimientos funcionales y no funcionales del
\nombreSistema, un sistema web de administración con control de acceso basado
en roles (RBAC). El documento está dirigido a:

\begin{itemize}
    \item El equipo docente y evaluador del curso \cursoDoc, como evidencia
          formal del alcance implementado.
    \item Los desarrolladores actuales y futuros del sistema, como referencia
          de mantenimiento y evolución.
    \item Cualquier interesado que requiera comprender el funcionamiento del
          sistema sin examinar su código fuente.
\end{itemize}

\section{Alcance}

El producto descrito se denomina \textbf{\nombreSistema} (nombre interno del
proyecto: \textit{Proyecto\_pw}). Es una aplicación web de tipo página única
(SPA) con arquitectura \textit{shell + frames} que provee:

\begin{itemize}
    \item Autenticación de usuarios con captcha propio, bloqueo por intentos
          fallidos y sesiones con token de expiración.
    \item Un sistema RBAC de dos niveles: acceso a pantallas (\textit{Frames})
          y acciones internas dentro de cada pantalla, configurable por rol
          desde la interfaz.
    \item Administración de usuarios, roles y permisos (exclusiva del
          Administrador).
    \item Un menú de navegación personalizable por cada usuario mediante
          arrastrar y soltar (SuperMenus propios), con disponibilidad global
          de opciones controlada por el Administrador.
    \item Un Dashboard con accesos rápidos personalizables, estadísticas
          (solo Administrador) y frases motivadoras (solo usuarios estándar).
    \item Gestión del perfil propio (consulta de datos y cambio de contraseña).
    \item Cinco Frames de ejemplo que demuestran el sistema de permisos,
          incluido un CRUD de tareas con permisos por operación.
\end{itemize}

Queda fuera del alcance del sistema: la recuperación de contraseñas por
correo, el registro público de usuarios (las cuentas solo las crea el
Administrador), la internacionalización (la interfaz es únicamente en
español) y la exposición de una API pública para terceros.

\section{Definiciones, acrónimos y abreviaturas}

\begin{longtable}{|>{\bfseries}p{3.6cm}|p{11cm}|}
\hline
\rowcolor{azulcorp!15} \textbf{Término} & \textbf{Definición} \\ \hline
\endhead
Frame & Pantalla del sistema que se carga dentro del marco de contenido del
        shell (por ejemplo: Usuarios, Dashboard, Frame 1). \\ \hline
ItemMenu & Opción de menú que abre un Frame. Es global al sistema y se
        registra en la tabla \texttt{menu}. \\ \hline
SuperMenu & Contenedor de ItemMenus creado por un usuario para organizar su
        menú personal. Es individual: cada usuario posee los suyos. \\ \hline
Shell & Página contenedora única (\texttt{dashboard/index.html}) que aloja la
        barra lateral, la barra superior y el marco donde se cargan los
        Frames sin recargar la estructura. \\ \hline
Módulo & Identificador RBAC asociado a un Frame (p.~ej. \texttt{usuarios},
        \texttt{frame1}), sobre el que se otorgan permisos. \\ \hline
Acción & Operación autorizable dentro de un módulo: \texttt{leer},
        \texttt{crear}, \texttt{editar} o \texttt{eliminar}. \\ \hline
Registro de Frames & Catálogo central del servidor
        (\texttt{servidor/permisos/registro.php}) que define los módulos
        existentes y sus acciones internas. \\ \hline
Semilla de menú & Estructura de menú inicial que se genera automáticamente
        una única vez al crear un usuario. \\ \hline
Token de sesión & Cadena hexadecimal de 64 caracteres generada
        aleatoriamente que identifica una sesión activa. \\ \hline
Captcha & Verificación gráfica generada en un lienzo (canvas) del navegador
        con seis caracteres aleatorios. \\ \hline
RBAC & \textit{Role-Based Access Control}: control de acceso basado en roles. \\ \hline
ERS & Especificación de Requerimientos de Software. \\ \hline
SPA & \textit{Single Page Application}: aplicación de página única. \\ \hline
CRUD & Crear, Leer, Actualizar y Eliminar (\textit{Create, Read, Update,
        Delete}). \\ \hline
PDO & \textit{PHP Data Objects}: capa de acceso a datos de PHP. \\ \hline
JSON & \textit{JavaScript Object Notation}: formato de intercambio de datos. \\ \hline
DnD & \textit{Drag \& Drop}: interacción de arrastrar y soltar. \\ \hline
bcrypt & Algoritmo de hash utilizado para almacenar contraseñas. \\ \hline
\end{longtable}

\section{Referencias}

\begin{itemize}
    \item IEEE Std 830-1998, \textit{IEEE Recommended Practice for Software
          Requirements Specifications}.
    \item Código fuente del proyecto \textit{Proyecto\_pw} (carpetas
          \texttt{paginas/}, \texttt{scripts/}, \texttt{servidor/},
          \texttt{estilos/}).
    \item Volcado de la base de datos \texttt{proyecto\_pw} (MySQL 8.0),
          archivo \texttt{proyecto\_pw.sql}.
    \item Documentación oficial de PHP 7.3 y MySQL 8.0.
\end{itemize}

\section{Visión general del documento}

El resto del documento se organiza de la siguiente manera: el
Capítulo~\ref{cap:descripcion} presenta la descripción general del producto,
su arquitectura, funciones, usuarios, restricciones y dependencias; el
Capítulo~\ref{cap:rf} especifica los requerimientos funcionales agrupados por
módulo; el Capítulo~\ref{cap:rnf} recoge los requerimientos no funcionales;
el Capítulo~\ref{cap:cu} desarrolla los casos de uso; el
Capítulo~\ref{cap:rn} enumera las reglas de negocio implementadas con su
evidencia en el código; finalmente, el Capítulo~\ref{cap:anexos} contiene los
anexos técnicos (modelo de datos, matriz de permisos, catálogo de endpoints y
estructura del proyecto).
